
\(2^\circ\)) \(1,\ldots,\xi^{r+1}\)
sont linéairement dépendants sur \(K^{\bullet}(S)\).

\(3^\circ\)) \(1,\ldots,\xi^r\) sont linéairement indépendants sur \(K^{\bullet}(S)\).

1.2.
Pour plus de commodité, nous allons rappeler ici quelques définitions et propriétés classiques des Modules gradués.

Soient \(S\) un schéma, \(\underline{B}\) une \(\underline{O}_S\)-Algèbre graduée, \(\underline{M}\) un \(\underline{B}\)-Module gradué, Pour tout \(\underline{B}\)-Module gradué \(\underline{N}\), on désigne comme d'habitude par \(\underline{N}(n)\) le \(\underline{B}\)-Module gradué défini par :
\[
  (\underline{N}(n))_k=\underline{N}_{n+k}\quad\text{pour tout }k.
\]
On dit que \(\underline{M}\) est un \(\underline{B}\)-Module gradué localement libre (resp. localement libre de type fini) si tout point de \(S\) possède un voisinage ouvert \(U\) tel que \(\underline{M}|U\) soit isomorphe, en tant que \(\underline{B}\)-Module gradué, à une somme directe de Modules de la forme \(\underline{B}(n)|U\) (resp. une somme directe finie).

On dit que \(\underline{M}\) est un \(\underline{B}\)-Module gradué de présentation finie si tout point de \(S\) possède un voisinage ouvert \(U\) tel qu'il existe une suite exacte \(\underline{P}\longrightarrow \underline{Q}\longrightarrow \underline{M}|U\longrightarrow 0\), où \(\underline{P}\) et \(\underline{Q}\) sont des \((\underline{B}|U)\)-Modules gradués libres de type fini, et les homomorphismes compatibles aux graduations.

1.2.1. Pour tout entier \(n\), on peut définir un homomorphisme canonique de \(\underline{B}\)-Modules gradués :
\begin{equation}\tag{1.2.1}
  \underline{M}_n\otimes_{\underline{O}_S} \underline{B}(-n)\longrightarrow M,
\end{equation}
en envoyant \(x\otimes a\), où \(x\) est une section de \(\underline{M}_n\), et \(a\) une section de \((\underline{B}(-n))_k=\underline{B}_{k-n}\) sur l'élément \(a.x\) de \(\underline{M}_k\).

1.2.2. On dit qu'un \(\underline{B}\)-Module gradué \(\underline{M}\) est libre de type fini, et homogène de degré \(n\) s'il est isomorphe en tant que \(\underline{B}\)-Module gradué à \((\underline{B}(-n))^k\) pour un certain entier \(k\).

